\chapter{Reliability, resilience, scalability, and capacity}

\section{Define reliable for whom}
Reliability is not “the server did not crash.” It is the probability that a
system performs its required function for a stated user, time, and environment.
An API can be available while returning stale or incorrect data; a batch job can
be correct while missing its deadline. Define the service indicator before
choosing a target.

Useful indicators include successful request ratio, freshness, correctness rate,
queue age, and latency percentiles. The percentile must include a population,
window, and boundary. “p95 under 200 ms” is incomplete unless it says which
requests, which route, whether retries are included, and where time starts and
ends.

\section{SLOs and error budgets}
A service-level objective (SLO) is a target such as “99.9 percent of eligible
requests in a rolling 30-day window succeed within the defined latency bound.”
The remaining error budget is the permitted failure fraction. It is a decision
tool: when the budget is healthy, a team may accept a risky improvement; when it
is exhausted, stability work has priority. Google describes SLOs as a way to
state and measure desired reliability \cite{google2016sre}.

The arithmetic is simple. For a 30-day window, 99.9 percent availability allows
about 43.2 minutes of unavailable time if the measurement is continuous. The
number is not a promise: maintenance, exclusions, sampling, and user impact
change the interpretation. State the measurement method.

\section{Failure modes and resilience}
Resilience is the ability to continue useful operation or recover after failure.
List failure modes by layer:

\begin{itemize}
  \item input: malformed, oversized, duplicated, or adversarial;
  \item process: crash, memory exhaustion, deadlock, resource leak;
  \item dependency: timeout, partial response, incompatible version;
  \item data: corruption, lag, missing migration, unrecoverable deletion;
  \item human and organizational: bad release, unclear ownership, unavailable expert.
\end{itemize}

For each, choose a response: prevent, detect, degrade, retry, isolate, fail
closed, fail open, or recover. “Fail gracefully” is not a response until the
degraded behaviour is named.

\section{Scalability and capacity}
Scalability is how resource use and performance change as workload grows. It is
not synonymous with speed. Capacity planning starts with a workload model:
arrival rate, request mix, payload size, concurrency, data growth, and burst
shape. Measure service time and a bottleneck resource. Little's Law gives a
useful relationship for a stable queueing system:

\[
  L = \lambda W,
\]

where $L$ is average work in the system, $\lambda$ is average arrival rate, and
$W$ is average time in the system. The relationship does not by itself predict
tail latency or behaviour during overload, but it exposes inconsistent capacity
claims.

\section{Load, stress, and performance tests}
A performance test is only meaningful with a stated workload, environment,
warm-up, data shape, concurrency, and measurement. A benchmark that reports
average latency may hide a failing tail. A stress test beyond capacity can reveal
recovery behaviour, but it should not be run against a production dependency
without explicit authorization and safeguards.

Optimization should follow a measured bottleneck. Cache design changes
correctness and invalidation; parallelism changes ordering and contention; a
database index changes write cost and storage. Record the baseline, the change,
and the new workload result.

\section{Backups and recovery}
A backup that has never been restored is an untested hope. Define recovery point
objective (RPO) and recovery time objective (RTO), encrypt backups, control
access, test restoration, and document dependencies such as keys and schema
versions. A replication system can copy corruption or accidental deletion, so it
is not a substitute for historical recovery points.

\begin{exercise}
Write an SLO for an order-creation endpoint, calculate its monthly error budget,
and name three events that should consume it. Then specify an RPO and RTO for
orders and justify them in terms of user and business impact.
\end{exercise}

\section{Chapter summary}
Reliability starts with a defined user-visible indicator. Use SLOs to make trade-
offs explicit, model failure modes, plan capacity from workload, test performance
with reproducible conditions, and treat restoration as a capability that must be
exercised.

